% ----------------------
\subsection{Section 89 (27 February 1833)}
% ----------------------
	\label{DC89}
	
	% -----
	\subsubsection{Historical Background}
	% -----
		\label{DC89-HB}
		
		% --
		\paragraph{JSP}
		% --
		
			``While no contemporaneous sources describing the circumstances under which JS dictated this 27 February 1833 revelation have been located, later accounts indicate that it was recorded in connection with the activities of the School of the Prophets. According to Brigham Young, heavy tobacco use---in the form of both smoking and chewing---among members of the school, combined with Emma Smith's and others' complaints about cleaning tobacco juice from the floor, led JS `to inquire of the Lord with regard to use of tobacco' and `to the conduct of the elders with this particular practice.' This revelation---composed largely of warnings and counsel regarding not only the use of tobacco, but also the consumption of various foods, `hot drinks,' wine, and `Strong drinks'---was the result of his inquiries. Known among church members as the `Word of Wisdom,' referring to the opening phrase of the text, the revelation was evidently recorded in JS's translating room in Newel K. Whitney's store. Zebedee Coltrin, who was present, recalled JS coming out of his translation room and reading the revelation to over twenty members of the school then in attendance. Joel Johnson added that the revelation was given in the evening.
			
			``At the time this revelation was dictated, the temperance movement, as well as other dietary reform movements, was beginning to factor more prominently in American culture. Initially fueled by concerns about the physiological effects of alcohol, calls for more moderate, or temperate, use of alcohol, and even for complete abstinence, had become increasingly identified with Christian reform movements by this time and had resulted in the formation of thousands of temperance societies throughout the United States. In Kirtland, Ohio, the Kirtland Temperance Society, whose founding constitution boasted 239 signatories, held its first annual meeting on 6 October 1830, shortly before Mormon missionaries arrived in the area from New York. The society continued for five years, during which time a distillery at Kirtland and two in Mentor, Ohio, evidently closed for want of patronage. According to one reminiscent account, the society disbanded after many of the temperance workers moved away as the Mormon population in the area grew, suggesting that Mormon involvement with the society was limited.

			``Arguments against the use of alcohol and other items, including tea, coffee, and tobacco, could be found in the religious, medical, and popular publications of the time, while arguments promoting the health value of other foods and drinks were also prevalent in the period's literature. In January and February 1833, JS himself was subscribing to the \textit{American Revivalist, and Rochester Observer}, an evangelical weekly that regularly published articles on temperance.

			``The contents of this revelation appear to have been available to many church members within months after its recording. Sidney Gilbert made a private copy probably sometime in the summer of 1833 in Missouri, while Wilford Woodruff copied the revelation in the back of his personal copy of the Book of Commandments, probably before September 1835---when the Doctrine and Covenants, which contained this revelation, became available. It was also printed as a broadsheet around January 1834 in Kirtland. Charges against church members in Ohio and Pennsylvania for disobeying some of the revelation's instruction surfaced in February 1834, within a year of the revelation's dictation, while local church conferences as far away as New York and Maine were referencing the revelation by summer 1834. Eber D. Howe published the revelation in his book \textit{Mormonism Unvailed} in 1834.

			``Among the members of the Church of Christ there was apparently some question as to what the revelation meant by `hot drinks,' prompting JS and Hyrum Smith, according to one reminiscent account, to explicitly identify coffee and tea at a meeting in Kirtland in July 1833 as the `hot drinks' prohibited by the revelation.%
				\footnote{%
					% \label{}%
					JSP note: ``Johnson, Notebook, [1]. Hyrum Smith made the same point nine years later in Nauvoo; in 1870, Brigham Young also identified the `hot drinks' mentioned in the revelation as tea and coffee. (`The Word of Wisdom,' \textit{Times and Seasons}, 1 June 1842, 3:800; Brigham Young, in \textit{Journal of Discourses}, 30 Oct. 1870, 13:277.)'' The reference for the Johnson notebook is ``Johnson, Joel H. Notebook, not before 1879. Joel Hills Johnson, Papers, ca. 1835--1882. CHL. MS 1546, fd. 7.''
					}
			Similarly, opinions on how strictly the revelation's instructions were to be followed appear to have differed among church members, probably as a result of the revelation's opening statement that the Word of Wisdom was given `for the benefit' of church members, `not by commandment or Constraint.' Possibly complicating the situation were the different ways this statement was presented. With one possible exception, the earliest manuscript versions of this revelation present the opening statement as part of the revelation, as do the printed broadsheet and Howe's \textit{Mormonism Unvailed}. In the 1835 edition of the Doctrine and Covenants, however, the opening paragraph appears as an italicized heading, allowing for the later interpretation that it was an introductory statement rather than part of the revelation proper and that, therefore, it was JS or one of his scribes, not God, who said the revelation was not a commandment.

			``In any event, the degree to which church members felt obligated to follow the revelation's instructions varied. Some, like Zebedee Coltrin, Joel Johnson, and John Tanner, chose to abstain from tea, coffee, tobacco, and alcohol almost immediately; others, like JS's wife Emma, who offered weary travelers tea and coffee upon their arrival in Kirtland in May 1833, apparently felt that using at least some of the items listed in the revelation was acceptable under some circumstances. A May 1835 letter from William W. Phelps to his wife, Sally Waterman Phelps---in which he spoke about the `sameness' of the Kirtland church members as they `drink cold water; and don't even mention tea and Coffee'---suggests the revelation was more universally understood among church members by that time, although many exceptions continue to appear in the historical record.

			``Many church members, for instance, apparently felt that it was acceptable for tea or alcohol to be taken medicinally. Emma may have offered tea and coffee to new arrivals with this idea in mind; JS's administering whiskey in June 1834 to George A. Smith, who was suffering from cholera, almost certainly reflected such an interpretation. In spite of JS's acquiescence with this practice, not everyone agreed with it. On 4 December 1836, for example, at the instigation of Sidney Rigdon, a meeting of church members in Kirtland voted unanimously to `discountenance the use intirely of all liquors from the Church in Sickness \& in health.' Over the ensuing years, nevertheless, various church members, including JS, continued to allow for the use of these drinks in cases of sickness.

			``JS and many others also allowed for a relaxed standard in adhering to the revelation's instructions during times of unusual difficulty and hardship. While overseeing a mass exodus of church members from Kirtland in spring 1838, for example, the seventies drew up a `constitution' charging leaders to see that `the word of wisdom [was] heeded'---that is, that `no tobacco, tea, coffee, snuff nor ardent spirits of any kind, [were] taken internally.' Hyrum Smith, however, speaking a few days later, advised those leaving Kirtland `not to be too particular in regard to the word of wisdom,' though subsequent events suggest his counsel was largely ignored. Following the collective trauma of the forced exodus from Missouri in winter 1838--1839, the prohibitions of the Word of Wisdom generally received less emphasis than they had earlier. Helen Mar Kimball Whitney, in a reminiscent account, reported that JS told church members suffering from malaria in early Nauvoo to `make tea and drink it' when the river water was unsuitable for drinking and that he `often made tea and administered it with his own hands.' Perhaps the best illustration of a more relaxed position regarding the Word of Wisdom in times of stress is John Taylor's account of events leading to JS's death at Carthage, Illinois, in June 1844, in which he noted that JS and his companions, who were feeling `unusually dull and languid' after several days of incarceration, drank some wine to raise their spirits.

			``In accordance with the revelation's provision that homemade wine could be taken when church members met `to offer up [their] sacrament' before God, church members continued to use wine, generally fermented, in the sacrament of the Lord's Supper. This same provision, coupled with the understanding that a sacrament was something `having a sacred character or function,' probably accounts for the times JS and others drank wine on several other occasions as well, including at the School of the Prophets, in various meetings held in the Kirtland temple, and at weddings. By the Nauvoo era, JS was more frequently making exceptions to the general observance of the Word of Wisdom that were not linked with health issues, hardship, or sacred functions, possibly indicating a more relaxed attitude on his part toward things like tea and wine. JS's journal for 11 March 1843, for example, indicates that he `had tea with his breakfast.' Two months later, on 3 May 1843, JS `drank a glass of wine with Sister Richards. of her Mothe[r]'s make, in England.' A year after that, JS `drank a glass of beer at Moissers [Frederick Moeser's].'

			``Only a few weeks prior to his death, JS seemed to reference the Word of Wisdom while counseling those who would be leaving to serve electioneering missions for his U.S. presidential campaign. At least in regard to alcohol, JS inveighed against drunkenness rather than just occasional consumption, which reflected his own actions in relation to alcohol. He informed the men that `we should never indulge our appetites to injure our influence, or wound the feelings of friends, or cause the spirit of the Lord to leave us. There is no excuse for any man to drink and get drunk in the church of Christ, or gratify any appetite, or lust, contrary to the principles of righteousness.' JS further instructed the men `on the principles of sobriety, and every thing pertaining to godliness at considerable length \& concluded by remarking that it is best to run on a long race and be careful to keep good wind \&c.'

			``Though the revelation instructed that meat was `to be used sparingly,' church members appear to have placed very little emphasis on that counsel, perhaps because this portion was rarely referenced by church leaders. Journals, reminiscences, and other personal records of the time that discuss specific provisions of the Word of Wisdom generally focus on the use of hot drinks, strong drinks, and tobacco rather than on the misuse or overuse of meat. The same is true of more official records and statements. In a noteworthy exception, Hyrum Smith, in an 1842 discourse on the Word of Wisdom, urged the Saints in Nauvoo to `attend to' the revelation's instructions regarding the use of meat and to be `sparing of the life of animals.'

			``Portions of this revelation reflect material in earlier JS revelations. Sometime around August 1830, a revelation on the emblems used in the sacrament of the Lord's Supper prohibited the members of the Church of Christ from purchasing wine and `strong drink' from their enemies and enjoined them not to partake of any such drink `except it is made new among you.' Another revelation, as well as the Book of Mormon, endorsed the use of herbs and other plants for treating the sick, and a 7 August 1831 revelation noted that animals, plants, and `all things which cometh of the earth in the season thereof is made for the benefit \& the use of man . . . to be used with judgement not to excess neither by extortion.' Still another revelation similarly observed that `the beasts of the field \& the fowls of the air \& that which cometh of the Earth is ordained for the use of man for food \& for raiment \& that he might have in abundance' and at the same time condemned anyone `that shedeth blood or that wasteth flesh \& hath no need.' This warning against wasting resources and food, especially meat, echoed JS's revision of Genesis 9:5 in which God tells Noah that `blood shall not be shed only for meat to save your lives and the blood of every beast will I require at your hands.'
			
			``The revelation makes no mention of an official penalty for disobeying its counsel, an issue that first presented itself to Kirtland leaders on 12 February 1834, when a `Bro Rich'---probably Leonard Rich---was `called in question for transgressing the word of wisdom.' A conference of church leaders and other men ordained to the priesthood forgave Rich `upon his promiseing to do better and reform his life.' The issue arose again eight days later---this time before the newly organized Kirtland high council---following a meeting that had been held in Pennsylvania in which `some of the members of that Church refused to partak[e] of the sacrament because the Elder administering it did not observe the words of wisdom to obey them.' Rather than addressing the Pennsylvanians' refusal, the Kirtland high council deliberated on the more fundamental issue of `whether disobedience to the word of wisdom was a transgression sufficient to deprive an official member from holding an office in the church, after haveing it sufficiently taught him.' The official decision, presented by JS and sanctioned by the council, was that `no official member in this church is worthy to hold an office after haveing the word of wisdom properly taught to him, and he, the official member, neglecting to comply with, or obey them.'

			``The council's decision was eventually published in the \textit{Messenger and Advocate} and appears to have been the basis for several policies and judgments made in Kirtland and Missouri. Records indicate that more severe actions, including excommunication, could be taken during this time when the violation of the principles taught in the revelation seemed particularly egregious or was part of a larger pattern of disobedience. Similarly, resolutions calling for the excommunication of church members who used `ardent spirits as a beverage' or who were `in the habit of drinking ardent spirits' were passed in various places in the early 1840s. At the same time, however, records from the Nauvoo period also indicate willingness on the part of church leaders and others to deal gently with those who were not obeying the revelation in the strictest sense and to give them time and reasons for reformation. Fearful that many church members were `following their old traditions,' for example, Hyrum Smith promised health, vigor, strength, and wisdom to those who kept the Word of Wisdom. An editorial in the \textit{Times and Seasons} counseled those who frequented `public places, where poison is dealt to the unwary' to be more actively engaged in the ministry to which they had been called, while those who used `tobacco and other intoxicating nauseates' were reminded that such substances `destroy the influence of the Holy Spirit.' Though disobedience to the Word of Wisdom was occasionally grounds for losing one's office during the Kirtland years, twenty-two men who apparently struggled to keep all of its provisions were ordained elders on 10 April 1843. Missionaries, similarly, were promised blessings if they kept the Word of Wisdom rather than being threatened with losing their licenses if they did not.

			``The copy of the revelation featured here is the private copy made by Sidney Gilbert. Several pieces of textual evidence, including the lack of clarifying and elaborating phrases that occur in other early copies, suggest that it may best represent the earliest version of the revelation. In the following transcript, significant textual differences are noted between this copy and the copy made by Frederick G. Williams in Revelation Book 2 in Kirtland, which was probably the earliest copy made in an official church record. All other early versions of this revelation closely follow the wording of Revelation Book 2.''
			
	% -----
	\subsubsection{Versions}
	% -----
		\label{DC89-V}
		
		See the \href{https://drive.google.com/file/d/1goAvNz8jS4R8slYfMG_db55Ty2z_vmgK/view?usp=sharing}{sections comparison} document.
	
		\begin{compactdesc}
			\item[JSP] \hyperref[EMS-1833-JS-27-FEB-1833]{27 February 1833}
			\item[BC] ---
			\item[DC 1835] Section 80, 207--208
			\item[MS] 1:9 (January 1841), 226--227
			\item[DC 1844] Section 81, 318--319
			\item[TS] 5:23 (15 December 1844), 736
		\end{compactdesc}

	% -----
	\subsubsection{Section 89:1} % *DC89-1 
	% -----
		\label{DC89-1}

		A WORD OF WISDOM, for the benefit of the council of high priests, assembled in Kirtland, and the church, and also the saints in Zion---

		% \begin{framed}
		% \scriptsize
			% \begin{compactdesc}
				% \item[JSP] 
				% % \item[BC] 
				% % \item[]
			% \end{compactdesc}
		% \end{framed}

	% -----
	\subsubsection{Section 89:2} % *DC89-2 
	% -----
		\label{DC89-2}

		To be sent greeting; not by commandment or constraint,%
			\footnote{%
				% \label{}%
				The Lord is sometimes explicit about whether what he's giving is a ``commandment''---see \hyperref[DC70-1]{DC 70:1--2}, for example.
				}
		but by revelation and the word of wisdom, showing forth the order and will of God in the temporal salvation%
			\footnote{%
				% \label{}%
				I think this phrase is the key to the whole section. The Word of Wisdom is about our \textit{temporal salvation} (for the phrase, see \hyperref[SALVATION-PHRASES]{here}). \hyperref[SALVATION-JS]{JS says} that ``salvation'' is to be saved from all our enemies, so a temporal salvation would be to be saved from all our temporal enemies. If you think about ill health and addictive behavior as enemies, then you can see why this revelation would be relevant to your temporal salvation.
				}
		of all saints in the last days---

		% \begin{framed}
		% \scriptsize
			% \begin{compactdesc}
				% \item[JSP] 
				% % \item[BC] 
				% % \item[]
			% \end{compactdesc}
		% \end{framed}

	% -----
	\subsubsection{Section 89:3} % *DC89-3 
	% -----
		\label{DC89-3}

		Given for a principle with promise, adapted to the capacity of the weak and the weakest of all saints, who are or can be called saints.

		% \begin{framed}
		% \scriptsize
			% \begin{compactdesc}
				% \item[JSP] 
				% % \item[BC] 
				% % \item[]
			% \end{compactdesc}
		% \end{framed}

	% -----
	\subsubsection{Section 89:4} % *DC89-4 
	% -----
		\label{DC89-4}

		Behold, verily, thus saith the Lord unto you: In consequence of evils and designs which do and will exist in the hearts of conspiring men in the last days, I have warned you, and forewarn you, by giving unto you this word of wisdom by revelation---

		% \begin{framed}
		% \scriptsize
			% \begin{compactdesc}
				% \item[JSP] 
				% % \item[BC] 
				% % \item[]
			% \end{compactdesc}
		% \end{framed}

	% -----
	\subsubsection{Section 89:5} % *DC89-5 
	% -----
		\label{DC89-5}

			\footnote{%
				% \label{}%
				I find this guidance remarkable, in light of how it's supposed to be for our ``temporal salvation.'' 
				}%
		That inasmuch as any man drinketh wine or strong drink among you, behold it is not good, neither meet in the sight of your Father, only in assembling yourselves together to offer up your sacraments before him.

		% \begin{framed}
		% \scriptsize
			% \begin{compactdesc}
				% \item[JSP] 
				% % \item[BC] 
				% % \item[]
			% \end{compactdesc}
		% \end{framed}

	% -----
	\subsubsection{Section 89:6} % *DC89-6 
	% -----
		\label{DC89-6}

		And, behold, this should be wine, yea, pure wine of the grape of the vine, of your own make.

		% \begin{framed}
		% \scriptsize
			% \begin{compactdesc}
				% \item[JSP] 
				% % \item[BC] 
				% % \item[]
			% \end{compactdesc}
		% \end{framed}

	% -----
	\subsubsection{Section 89:7} % *DC89-7 
	% -----
		\label{DC89-7}

		And, again, strong drinks are not for the belly, but for the washing of your bodies.

		% \begin{framed}
		% \scriptsize
			% \begin{compactdesc}
				% \item[JSP] 
				% % \item[BC] 
				% % \item[]
			% \end{compactdesc}
		% \end{framed}

	% -----
	\subsubsection{Section 89:8} % *DC89-8 
	% -----
		\label{DC89-8}

			\footnote{%
				% \label{}%
				This is maybe the most amazing part of the whole revelation, to prohibit tobacco.
				}%
		And again, tobacco is not for the body, neither for the belly, and is not good for man, but is an herb for bruises and all sick cattle, to be used with judgment and skill.

		% \begin{framed}
		% \scriptsize
			% \begin{compactdesc}
				% \item[JSP] 
				% % \item[BC] 
				% % \item[]
			% \end{compactdesc}
		% \end{framed}

	% -----
	\subsubsection{Section 89:9} % *DC89-9 
	% -----
		\label{DC89-9}

		And again, hot drinks are not for the body or belly.

		% \begin{framed}
		% \scriptsize
			% \begin{compactdesc}
				% \item[JSP] 
				% % \item[BC] 
				% % \item[]
			% \end{compactdesc}
		% \end{framed}

	% -----
	\subsubsection{Section 89:10} % *DC89-10 
	% -----
		\label{DC89-10}

		And again, verily I say unto you, all wholesome herbs God hath ordained for the constitution, nature, and use of man---

		% \begin{framed}
		% \scriptsize
			% \begin{compactdesc}
				% \item[JSP] 
				% % \item[BC] 
				% % \item[]
			% \end{compactdesc}
		% \end{framed}

	% -----
	\subsubsection{Section 89:11} % *DC89-11 
	% -----
		\label{DC89-11}

		Every herb in the season thereof, and every fruit in the season thereof; all these to be used with prudence and thanksgiving.

		% \begin{framed}
		% \scriptsize
			% \begin{compactdesc}
				% \item[JSP] 
				% % \item[BC] 
				% % \item[]
			% \end{compactdesc}
		% \end{framed}

	% -----
	\subsubsection{Section 89:12} % *DC89-12 
	% -----
		\label{DC89-12}

		Yea, flesh also of beasts and of the fowls of the air, I, the Lord, have ordained for the use of man with thanksgiving; nevertheless they are to be used sparingly;%
			\footnote{%
				% \label{}%
				Note the next few verses and compare \hyperref[DC49-18]{DC 49:18--19}.
				}

		% \begin{framed}
		% \scriptsize
			% \begin{compactdesc}
				% \item[JSP] 
				% % \item[BC] 
				% % \item[]
			% \end{compactdesc}
		% \end{framed}

	% -----
	\subsubsection{Section 89:13} % *DC89-13 
	% -----
		\label{DC89-13}

		And it is pleasing unto me that they should not be used, only in times of winter, or of cold, or famine.

		% \begin{framed}
		% \scriptsize
			% \begin{compactdesc}
				% \item[JSP] 
				% % \item[BC] 
				% % \item[]
			% \end{compactdesc}
		% \end{framed}

	% -----
	\subsubsection{Section 89:14} % *DC89-14 
	% -----
		\label{DC89-14}

		All grain is ordained for the use of man and of beasts, to be the staff of life, not only for man but for the beasts of the field, and the fowls of heaven, and all wild animals that run or creep on the earth;

		% \begin{framed}
		% \scriptsize
			% \begin{compactdesc}
				% \item[JSP] 
				% % \item[BC] 
				% % \item[]
			% \end{compactdesc}
		% \end{framed}

	% -----
	\subsubsection{Section 89:15} % *DC89-15 
	% -----
		\label{DC89-15}

		And these hath God made for the use of man only in times of famine and excess of hunger.

		% \begin{framed}
		% \scriptsize
			% \begin{compactdesc}
				% \item[JSP] 
				% % \item[BC] 
				% % \item[]
			% \end{compactdesc}
		% \end{framed}

	% -----
	\subsubsection{Section 89:16} % *DC89-16 
	% -----
		\label{DC89-16}

		All grain is good for the food of man; as also the fruit of the vine; that which yieldeth fruit, whether in the ground or above the ground---

		% \begin{framed}
		% \scriptsize
			% \begin{compactdesc}
				% \item[JSP] 
				% % \item[BC] 
				% % \item[]
			% \end{compactdesc}
		% \end{framed}

	% -----
	\subsubsection{Section 89:17} % *DC89-17 
	% -----
		\label{DC89-17}

		Nevertheless, wheat for man, and corn for the ox, and oats for the horse, and rye for the fowls and for swine, and for all beasts of the field, and barley for all useful animals, and for mild drinks, as also other grain.

		% \begin{framed}
		% \scriptsize
			% \begin{compactdesc}
				% \item[JSP] 
				% % \item[BC] 
				% % \item[]
			% \end{compactdesc}
		% \end{framed}

	% -----
	\subsubsection{Section 89:18} % *DC89-18 
	% -----
		\label{DC89-18}

		And all saints who remember to keep and do these sayings, walking in obedience to the commandments, shall receive health in their navel and marrow to their bones;%
			\footnote{%
				% \label{}%
				This last part of the verse reminds of some of the things said at the veil in the \hyperref[TEMPLE-ENDOWMENT-ENDOW-PRE-1990-PASSING]{endowment ceremony}.
				} 

		% \begin{framed}
		% \scriptsize
			% \begin{compactdesc}
				% \item[JSP] 
				% % \item[BC] 
				% % \item[]
			% \end{compactdesc}
		% \end{framed}

	% -----
	\subsubsection{Section 89:19} % *DC89-19 
	% -----
		\label{DC89-19}

		And shall find wisdom and great treasures of knowledge, even hidden treasures;

		% \begin{framed}
		% \scriptsize
			% \begin{compactdesc}
				% \item[JSP] 
				% % \item[BC] 
				% % \item[]
			% \end{compactdesc}
		% \end{framed}

	% -----
	\subsubsection{Section 89:20} % *DC89-20 
	% -----
		\label{DC89-20}

		And shall run and not be weary, and shall walk and not faint.

		% \begin{framed}
		% \scriptsize
			% \begin{compactdesc}
				% \item[JSP] 
				% % \item[BC] 
				% % \item[]
			% \end{compactdesc}
		% \end{framed}

	% -----
	\subsubsection{Section 89:21} % *DC89-21 
	% -----
		\label{DC89-21}

		And I, the Lord, give unto them a promise, that the destroying angel shall pass by them,%
			\footnote{%
				% \label{}%
				Pointing back up to the temporal salvation above, in \hyperref[DC89-2]{verse 2}?
				}
		as the children of Israel, and not slay them.  Amen.

		% \begin{framed}
		% \scriptsize
			% \begin{compactdesc}
				% \item[JSP] 
				% % \item[BC] 
				% % \item[]
			% \end{compactdesc}
		% \end{framed}